%% ch01 --- CPython, the GIL and the free-threaded build
%%
%% Stub, generated by tools/gen_stubs.py from tools/chapters.json.
%% The brief below is the contract for this chapter: what it must argue, what
%% it must buy the reader, and what it must NOT repeat from the companion
%% volumes. Writing the chapter means deleting the stub block below (the one
%% that opens on the next \chapter line) and replacing it with the chapter
%% text -- after which this file is never regenerated.
%% If the brief turns out to be wrong once the code has been run, change it in
%% the manifest and record the contradiction in CLAUDE.md under *Resolved
%% questions*.

\chapter{CPython, the GIL and the free-threaded build}
\label{ch:runtime}

\chapterstub{%
Argument: the assumptions a CLR engineer carries about the runtime stop
holding one at a time, and each one has a name. There is no JIT you can
count on by default (the one in 3.14 is experimental and off); one process
is one interpreter with one global interpreter lock that serialises Python
bytecode across threads; and the free-threaded build (PEP 703, officially
supported since 3.14 under PEP 779) is a separate binary,
\code{python3.14t}, rather than a flag. Bytecode and \code{\_\_pycache\_\_}
against IL and the JIT; \code{python -m} against \code{dotnet run}; a module
against an assembly; \code{sys.path} against assembly probing; what \code{if
\_\_name\_\_ == "\_\_main\_\_"} is actually for. Payoff: the reader can say
which of four workloads the GIL hurts and measure it \dash{} experiment E1,
CPU-bound work on the default and the free-threaded build, four threads, on
the pinned interpreter. Traps to elicit: \enquote{threads are useless in
Python} (false for I/O), \enquote{the free-threaded build is the default
now} (false), \enquote{Python has no compile step} (bytecode is compiled and
cached). Listings: \code{python -m dis} on a five-line function; a thread
pool on CPU-bound and on I/O-bound work, timed; detecting the free-threaded
build from \code{sys}. Three exercises. Overlap rule: the LangChain book's
chapter 3 covers the GIL in the context of an event loop; this chapter owns
the runtime model and points there, in one sentence, for the async
consequence.

\medskip
\textbf{Word budget:} 3000.\\
\textbf{Ships in:} v0.1.\\
\textbf{Measurement:} E1.
}
